\section{Exchange of Derivative with Limit}

\begin{theorem}[Exchange of Limit with Derivative]
\mymark{***}
Let $I\subset \RR$ be an interval and let $f_k\in C^1(I)$ be functions such that
$f_k(x_0)$ converges for at least one point $x_0\in I$ and the sequence of
derivatives $f_k'$ converges to a function $g\colon I \to \RR$ uniformly on
every closed and bounded interval $[a,b]\subset I$. Then there exists $f\in
C^1(I)$ such that $f'=g$ and $f_k$ converges to $f$ uniformly on every closed
and bounded interval $[a,b]\subset I$. Under these hypotheses, the derivative
can be exchanged with the limit:
\[
  \lim_{k\to +\infty}\enclose{\frac{d}{dx} f_k(x)}
  = f'(x)
  = \frac{d}{dx} \enclose{\lim_{k\to +\infty} f_k(x)},
  \qquad \forall x \in I.
\]
\end{theorem}
%
\begin{proof}
\mymark{***}
By hypothesis, there exists $y_0\in \RR$ such that $f_k(x_0) \to y_0$. Define
\[
  f(x) = y_0 + \int_{x_0}^x g(t)\, dt.
\]
By the continuity of the uniform limit, we know that $g$ is continuous, thus we
can apply the fundamental theorem of calculus to deduce that $f'=g$. We now show
that on every interval $[a,b]\subset I$, we have $f_k \To f$. By the fundamental
theorem of integral calculus, we have:
\[
  \int_{x_0}^x f_k'(t) dt = f_k(x) - f_k(x_0)
\]
thus
\begin{align*}
  \sup_{x\in [a,b]}\abs{f_k(x) - f(x)}
    &= \sup_{x\in [a,b]} \abs{f_k(x_0) + \int_{x_0}^x f_k'(t) - y_0 -
  \int_{x_0}^x g(t)\, dt} \\
    &\le \abs{f_k(x_0) - y_0} + \sup_{x\in [a,b]}\abs{\int_{x_0}^x \abs{f_k'(t)
  - g(t)}\, dt} \\
  &\le \abs{f_k(x_0) - y_0} + (b-a)\Abs{f_k' - g} \to 0.
\end{align*}
\end{proof}

The space $C^1([a,b])$ is a vector subspace of $C^0([a,b])$ but it is not
closed, as deduced from example~\ref{ex:466533} (one could even prove that $C^1$
is dense in $C^0$), thus $C^1$ is not complete with respect to the uniform norm.
To transform the space $C^1([a,b])$ into a Banach space, we can define a
stronger norm, such as the following:
\[
  \Abs{f}_{C^1} = \Abs{f}_\infty + \Abs{f'}_\infty.
\]
\begin{theorem}[$C^1$ Banach Space]
\mymark{*}
The vector space $C^1([a,b])$ endowed with the norm $\Abs{\cdot}_{C^1}$ is a
Banach space.
\end{theorem}
%
\begin{proof}
It is easy to verify that $\Abs{\cdot}_{C^1}$ is a norm on $C^1([a,b])$, we only
need to verify that the space is complete. Let $f_k$ be a Cauchy sequence with
respect to the $C^1$ norm. Then $f_k'$ and $f_k$ are both Cauchy sequences in
$C^0$ since $\Abs{f_k}_\infty \le \Abs{f_k}_{C^1}$ and $\Abs{f_k'}_\infty \le
\Abs{f_k}_{C^1}$. Therefore, by the completeness of $C^0$, we know that there
exist $f,g\in C^0([a,b])$ such that $f_k\To f$ and $f_k'\To g$. Based on the
theorem of exchange of limit with derivative, we can assert that $f\in C^1$ and
$f'=g$, thus
\[
  \Abs{f_k-f}_{C^1} = \Abs{f_k-f}_\infty + \Abs{f_k'-g}_\infty \to 0.
\]
\end{proof}

The theorem of exchange of limit with integral tells us that the integral
operator $S\colon C^0 \to C^1$ is continuous between the two Banach spaces. The
differential operator $D\colon C^1 \to C^0$ $f\mapsto Df = f'$ is obviously
continuous as well.

\section{Series of Functions}

If $f_k\colon A \to \RR$ is a sequence of functions defined on the same set $A$,
we can consider (as we have already done for numerical sequences) the sequence
of partial sums:
\[
  S_n(x) = \sum_{k=0}^n f_k(x), \qquad x\in A.
\]
This sequence is called the \emph{series} corresponding to the sequence of
functions $f_k$ and is sometimes denoted as $\sum f_n$. For each $x$ where the
series is convergent, we can define the \emph{sum}%
\mymargin{sum}\index{sum} of the series
\[
  S(x) = \sum_{k=0}^{+\infty} f_k(x) = \lim_{n\to +\infty} S_n(x).
\]
The sum $S$ is thus the pointwise limit of the sequence of partial sums $S_n$.

The theorems we have proven for sequences of functions are therefore valid for
series of functions as well. It suffices to remember that the \emph{uniform
convergence of the series}%
\mymargin{uniform convergence of a series}%
\index{series!uniform convergence}%
\index{convergence!uniform of a series}%
  is the uniform convergence of the partial sums. Thus $\sum f_k$ converges
 uniformly to $S$ if $S_n \To S$, that is, if
\[
  \Abs{S - S_n}_\infty = \Abs{\sum_{k=n+1}^{+\infty} f_k}_\infty \to 0
  \qquad \text{as $n\to +\infty$.}
\]

\begin{theorem}[Integral of a Series of Functions]
\label{th:integrale_serie}
\mymark{**}%
\index{theorem!integration of a series of functions}%
\index{series!integral}%
\mymargin{integral of a series}%
\index{integral!of a series}%
Let $f_k\colon [a,b]\to\RR$ be a sequence of continuous functions defined on the
interval $[a,b]\subset \RR$. If the series $\sum f_k$ converges uniformly, then
the integral can be exchanged with the sum of the series:
\[
  \int_a^b \enclose{\sum_{k=0}^{+\infty} f_k(t)}\, dt
  = \sum_{k=0}^{+\infty} \enclose{\int_a^b f_k(t)\, dt}
  \qquad \forall x \in I.
\]
\end{theorem}
\begin{proof}
\mymark{**}
The proof is a simple consequence of the fact that the exchange can be made on
finite sums and the passage to the limit can be made thanks to the theorem of
exchange of limit with integral.

Let $S_n = \sum f_n$ be the sequence of partial sums and let $S$ be the limit of
the partial sums. By hypothesis, $S_n\To S$. Applying the theorem of exchange of
integral with limit, we have
\[
  \lim_{n\to +\infty} \int_a^b S_n(t)\, dt = \int_a^b S(t)\, dt.
\]
But on one hand, exploiting the additivity of the integral on finite sums:
\begin{align*}
  \lim_{n\to +\infty} \int_a^b S_n(t)\, dt
   &= \lim_{n\to+\infty}\int_a^b \enclose{\sum_{k=0}^n f_k(t)} \,  dt\\
   &= \lim_{n\to+\infty}\sum_{k=0}^n \enclose{\int_a^b f_k(t)\, dt}\\
   &= \sum_{k=0}^\infty \enclose{\int_a^b f_k(t)\, dt}
\end{align*}
and on the other hand:
\[
  \int_a^b S(t)\, dt = \int_a^b \enclose{\sum_{k=0}^{+\infty} f_k(t)}\, dt.
\]
\end{proof}

\begin{theorem}[Derivative of a Series of Functions]
\mymark{**}
\index{theorem!differentiation of a series of functions}
\index{series!derivative}
\mymargin{differentiation of a series}
Let $f_k\colon I\to\RR$ be a sequence of continuous functions defined on the
interval $I$. If the functions $f_k$ are of class $C^1$ and the series of
derivatives $\sum f_k'$ converges uniformly on every closed and bounded interval
$[a,b]\subset I$ and if there is at least one point $x_0\in I$ such that the
series $\sum f_k(x_0)$ converges, then
\[
  \frac{d}{dx} \sum_{k=0}^{+\infty} f_k(x) = \sum_{k=0}^{+\infty} \frac{d}{dx}f_k(x)
  \qquad \forall x \in I.
\]
\end{theorem}

\begin{proof}
\mymark{**}
Let $S_n$ be the sequence of partial sums. By hypothesis, we know that there
exists a function $T\colon I \to \RR$ such that $S_n' \To T$ in every interval
$[a,b]\subset I$. We also know that $S_n(x_0)$ converges. Thus we can apply the
theorem of exchange of limit with derivative to obtain that there exists $S\in
C^1(I)$ such that $S_n(x)\to S(x)$ for every $x\in I$ and
\[
   S'(x) = T(x) \qquad \forall x\in I.
\]
But on one hand
\begin{align*}
S'(x)
&= \frac{d}{dx} \lim_{n\to +\infty} S_n(x) \\
&= \frac{d}{dx} \sum_{k=0}^{+\infty} f_k(x)
\end{align*}
and on the other hand
\begin{align*}
T(x)
&= \lim_{n\to +\infty} S_n'(x)
 = \lim_{n\to +\infty} \frac{d}{dx} \sum_{k=0}^n f_k(x) \\
&= \lim_{n\to +\infty} \sum_{k=0}^n f_k'(x)
 = \sum_{k=0}^{+\infty} f_k'(x).
\end{align*}
\end{proof}

The uniform convergence of a series is not very simple to verify. The following
condition is simpler and will be shown to be stronger.

\begin{definition}[Total Convergence of a Series of Functions]
\mymark{***}
Let $f_k\colon A \to \RR$ be functions defined on a set $A\subset \RR$. We say
that the series of functions $\sum f_k$ \emph{converges totally}%
\mymargin{total convergence}%
\index{convergence!total}%
if the numerical series $\sum \Abs{f_k}_\infty$ is convergent.
\end{definition}

\begin{theorem}[Total Convergence]
\mymark{***}
If the series $\sum f_n$ converges totally, then it converges uniformly.
\end{theorem}
%
\begin{proof}
\mymark{***}
For fixed $x$, the series $\sum f_n(x)$ converges absolutely since
\[
  \sum_{k=0}^\infty \abs{f_n(x)}
  \le \sum_{k=0}^\infty \Abs{f_n}_\infty < +\infty.
\]
Thus the series converges and setting
\[
  S_n(x) = \sum_{k=0}^n f_k(x), \qquad
  S(x) = \sum_{k=0}^{+\infty} f_k(x)
\]
we have that $S_n\to S$ pointwise. To show that $S_n \To S$, it suffices to
observe that as $n\to +\infty$, we have:
\[
  \abs{S(x) - S_n(x)}
  = \abs{\sum_{k=n+1}^{+\infty} f_k(x)}
  \le \sum_{k=n+1}^{+\infty}\abs{f_k(x)}
  \le \sum_{k=n+1}^{+\infty}\Abs{f_k}_\infty \to 0.
\]
\end{proof}

\begin{theorem}[Total Convergence of Power Series]
\label{th:convergenza_totale}
\mymark{***}%
Let $\sum a_n z^n$ be a power series and let $R\in[0,+\infty]$ be its radius of
convergence. Then the series converges totally on every disk $D_r$ with $r<R$.
\end{theorem}
%
\begin{proof}
Now observe that on the disk of radius $r$, we have $\Abs{a_k z^k}_\infty =
\abs{a_k} r^k$ and thus
\[
\sum_{k=0}^{+\infty} \Abs{a_k z_k}_\infty
= \sum_{k=0}^{+\infty} \abs{a_k}r^k < +\infty
\]
since the series $\sum a_k z^k$ converges absolutely for $z=r$ being $r<R$.
\end{proof}

\begin{corollary}
\label{cor:derivata_serie_potenze}
\mymark{**}
The power series
\[
  f(x) = \sum_{k=0}^{+\infty} a_k x^k
\]
has the same radius of convergence $R$ as the series of derivatives
\[
  g(x) = \sum_{k=1}^{+\infty} k a_k x^{k-1}
\]
and for $x\in (-R,R)$, we have
\[
  f'(x) = g(x).
\]
\end{corollary}
%
\begin{proof}
That the two series have the same radius of convergence has already been proven
in Theorem~\ref{th:raggio_serie_derivate}. In the previous theorem, we showed
that on every interval $[-r,r]$ with $r<R$, the power series with sum $f$
converges totally. Thus it converges uniformly, and we can exchange the
derivative with the sum to obtain $f'(x) = g(x)$.
\end{proof}

\begin{example}
We know that the power series
\[
f(x) = \sum_{k=1}^{+\infty} \frac{x^k}{k}
\]
has a radius of convergence $R=1$ (using, for example, the ratio test). The
series of derivatives is
\[
 g(x) = \sum_{k=1}^{+\infty} x^{k-1} = \sum_{k=0}^{+\infty}x^k = \frac{1}{1-x}.
\]
Thus for $\abs{x}<1$, we have
\[
  f'(x) = g(x) = \frac{1}{1-x}
\]
from which
\[
 f(x) = f(0) + \int_0^x f'(t)\, dt = \int_0^x \frac{1}{1-t}\, dt
  = \Enclose{-\ln(1-t)}_0^x = -\ln(1-x).
\]
We have thus discovered that
\[
  \sum_{k=1}^{+\infty} \frac{x^k}{k} = - \ln (1-x) \qquad\text{for every $x\in (-1,1)$.}
\]
Observe now that the series with sum $f(x)$ does not converge for $x=1$
(harmonic series) but converges for $x=-1$ (Leibniz's test). By
Theorem~\ref{th:lemma_abel} (Abel's lemma), we know that the function $f(x)$ is
continuous at the point $x=-1$, and thus we can conclude that
\[
  \sum_{k=1}^{+\infty} \frac{x^k}{k} = -\ln(1-x)
  \qquad \forall x \in [-1,1).
\]
In particular, we have found the sum of the alternating harmonic series:
\[
  \sum_{k=1}^{+\infty} \frac{(-1)^{k+1}}{k} = - f(x) = \ln 2.
\]

Observe that this information is consistent with the Taylor expansion of
$\ln(1+x)$ that we had already determined. However, it is not a consequence of
it, as the Taylor expansion gives us information only for $x\to 0$, whereas now
we have obtained information for every $x$ in $[-1,1)$.
 \end{example}

\begin{example}
We apply the previous idea to the function $\arctg$. We have
\[
  \arctg'(x) = \frac{1}{1+x^2} = \sum_{k=0}^{+\infty} (-x^2)^k
  = \sum_{k=0}^{+\infty} (-1)^k x^{2k}
  = \sum_{k=0}^{+\infty} \frac{(-1)^k}{2k+1}(x^{2k+1})'.
\]
The series
\[
 f(x) = \sum_{k=0}^{+\infty}\frac{(-1)^k}{2k+1} x^{2k+1}
\]
has a radius of convergence $R=1$, and thus for every $x\in(-1,1)$, we know that
the series of derivatives converges to the derivative of the series, from which
\[
  f'(x) = \arctg' x.
\]
Since $f(0) = 0 = \arctg 0$, we can conclude that $f(x) =\arctg x$ for every
$x\in (-1,1)$. The series is also convergent for $x=1$, by Leibniz's test. By
continuity (Abel's Lemma), we obtain that $f(1) = \arctg 1$. Thus
\[
  \arctg x = \sum_{k=0}^{+\infty}\frac{(-1)^k}{2k+1}x^{2k+1}
  \qquad \forall x \in (-1,1].
\]
In particular, for $x=1$, we obtain the \emph{Gregory-Leibniz} formula%
\mymargin{Gregory-Leibniz}\index{Gregory-Leibniz}
\index{$\pi$!Gregory-Leibniz formula}
\index{Gregory!approximation $\pi$}
\index{Leibniz!approximation $\pi$}
\index{formula!Gregory-Leibniz for $\pi/4$}
\[
  \frac{\pi}{4} = \sum_{k=0}^{+\infty} \frac{(-1)^k}{2k+1} =
   1 - \frac{1}{3} + \frac{1}{5} - \frac{1}{7} + \dots
\]
\end{example}

The following example uses total convergence to construct a continuous function
that is not differentiable at any point. Functions of this kind are not merely
mathematical curiosities but can be useful for modeling natural phenomena, such
as Brownian motion or the behavior of a stock index.

\begin{example}[Weierstrass Function]
Let $a=\frac 1 8$ and $b=512$, consider the function $W\colon \RR \to \RR$,
defined by
\[
   W(x) = \sum_{k=0}^{+\infty} a^k \cos(b^k \pi x).
\]
The function $W$ is continuous at every point of $\RR$ but is not differentiable
at any point of $\RR$.
\end{example}
\begin{proof}
Let $f_k(x) = a^k \cos(b^k \pi x)$, we have $\Abs{f_k}_{\infty} \leq a^k$, and
therefore, since $\sum a^k$ is convergent, the series of functions $W = \sum
f_k$ converges totally on $\RR$ and thus converges uniformly. Since each $f_k$
is continuous, $W$ is also continuous.

It is much more delicate to prove that $W$ is not differentiable at any point of
$\RR$. The idea is that the function $f_k$ has a derivative that oscillates with
amplitude $\pi a^k b^k$, while its modulus, as already seen, has amplitude
$a^k$. If we take two points that have a distance of the order of $b^m$, the
function $f_m$ is the one that contributes most to the difference of the
incremental ratio of $f$, since for $k<m$, the functions oscillate relatively
little, while for $k>m$, the functions have very small amplitude. We develop
this idea.

Fix any $x\in \RR$. Our goal is to find a sequence $\delta_n\to 0$ such that
\begin{equation}\label{eq:30234354}
  \abs{\frac{W(x+\delta_n) - W(x)}{\delta_n}} \to +\infty.
\end{equation}
This implies that $W$ is not differentiable at $x$. Since the cosine has period
$2\pi$, we are certain that in the interval $[4\pi,6\pi]$, the function $\theta
\mapsto \cos(b^n \pi x+\theta)$ assumes all values between $-1$ and $1$. Thus
there exists $\theta_n \in [4\pi,6\pi]$ such that
\[
   \abs{\cos (b^n \pi x + \theta_n) - \cos(b^n\pi x)} = 1.
\]
Set $\delta_n = \frac{\theta_n}{b^n \pi}$ so that
\[
  \abs{\cos (b^n \pi (x+\delta_n))- \cos(b^n \pi x)} = 1.
\]
We then use the following estimates
\[
  \abs{\cos(b^k\pi (x + \delta_n)) - \cos(b^k \pi x)}
  \begin{cases}
    \le b^k\pi\delta_n &\text{if $k<n$},\\
    = 1 &\text{if $k=n$},\\
    \le 2 &\text{if $k>n$}.
  \end{cases}
\]
We have therefore:
\begin{align*}
    \abs{\frac{W(x+\delta_n) - W(x)}{\delta_n}} 
        &= \frac{\pi b^n}{\theta_n}\cdot
    \abs{\sum_{k=0}^{+\infty}a^k\Enclose{\cos(b^k\pi (x + \delta_n) - \cos(b^k
    \pi x)}}  \\
        &\ge \frac{\pi b^n}{\theta_n}\cdot \Enclose{a^n - \sum_{k=0}^{n-1} a^k
    b^{k-n} \theta_n - \sum_{k=n+1}^{+\infty} 2a^k} \\
        &\ge \frac{a^n b^n}{6} -
    \frac{b^n}{4}\Enclose{\frac{6\pi}{b^n}\sum_{k=0}^{n-1} (ab)^k + 2 a^{n+1}
    \frac{1}{1-a}} \\
        &\ge \frac{a^n b^n}{6} -\frac {3\pi} 2\cdot  \frac{(ab)^n-1}{ab-1} -
    \frac{a^n b^n}{2}\cdot \frac{a}{a-1}\\
        &\ge (ab)^n \Enclose{\frac 1 6 - 5\cdot \frac{1}{ab-1} - \frac 1 2 \cdot
    \frac{1}{1-a}} \\
    &= (ab)^n \frac{(ab-1)(1-a)-30(1-a)-3a(ab-1)}{6(ab-1)(1-a)}.
\end{align*}
The choices of $a$ and $b$ ensure that we have $ab>1$, $1-a>0$, and
\begin{align*}
(ab-1)(1-a)-30(1-a)-3a(ab-1)
&= ab(1-4a) - 31 + 34a > 0.
\end{align*}
Thus, we obtain \eqref{eq:30234354} as desired.
\end{proof}